\clearpage
\section{Basic Topology}
\subsection*{Exercises}
\addcontentsline{toc}{subsection}{Exercises}

\begin{exercise}
Let $E'$ be the set of all limit points of a set $E$. 
Prove that $E'$ is closed.
Prove that $E$ and $\overline{E}$ have the same limit points. (Recall that $\overline{E}=E \cup E'$.)
Do $E$ and $E'$ always have the same limit points?
\end{exercise}

\begin{proof}
Let $p$ be a limit point of $E'$ and let $N$ be any neighborhood of $p$.
Then $N$ contains a point $q \neq p$ such that $q \in E'$.
Since $N$ is open, there exists a neighborhood $N_q$ of $q$ such that $N_q \subset N$.
Without loss of generality, assume that $N_q$ does not contain $p$.
For if $p \in N_q$ then we take $r=d(p,q)>0$, hence $p \notin N_r(q) \subset N_q \subset N$.

Recall that $q$ is a limit point of $E$ ($q \in E'$), we see that $N_q$ contains a point $s \neq q$ such that $s \in E$.
since $p \notin N_q$, we have $p \neq s$.
Therefore, $N \supset N_q$ contains a point $s \neq p$ such that $s \in E$.
Since $N$ was arbitrary, $p$ is a limit point of $E$ and $p \in E'$.

Since the choice of $p$ was arbitrary, every limit point of $E'$ is a point of $E'$.
Thus, $E'$ is closed.

Next, let us prove that $E$ and $\overline{E}$ have the same limit points.
Suppose $p$ is a limit point of $E$.
Then every neighborhood of $p$ contains a point $q \neq p$ such that $q \in E \subset E \cup E' = \overline{E}$.
Thus $p$ is a limit point of $\overline{E}$. 

Conversely, suppose $p$ is a limit point of $\overline{E}$. 
Then every neighborhood $N$ of $p$ contains a point $q \neq p$ such that $q \in \overline{E}$.
Note that $q \in E$ or $q \in E'$. 

If $q \in E$ then we are done. 
Otherwise, suppose $q \in E'$. 
Since $N$ is open, there is a neighborhood $N_q$ of $q$ such that $N_q \subset N$.
Without loss of generality, assume $N_q$ does not contain $p$.
Since $q$ is a limit point of $E$, $N_q$ contains a point $s \neq q$ such that $s \in E$.
Hence, $N$ contains a point $s \neq p$ (since $p \notin N_q$) such that $s \in E$.
Since $N$ was arbitrary, $p$ is a limit point of $E$.

We conclude that $E$ and $\overline{E}$ have the same limit points.

Finally, let $E=\left\{\frac{1}{n} \mid n \in \N\right\}$. It is clear that $E' = \{0\}$ but there is no limit point of $E'$. Hence it is possible that $E$ and $E'$ have different limit points.
\end{proof}
\begin{exercise}
Let $A_1, A_2, A_3, \ldots$ be subsets of a metric space.
\begin{enumerate}
    \item[(a)] If $B_n = \bigcup_{i=1}^nA_i$, prove that $\overline{B_n} = \bigcup_{i=1}^n \overline{A_i}$, for $n= 1, 2, 3, \ldots$.
    \item[(b)] If $B= \bigcup_{i=1}^{\infty} A_i$, prove that $\overline{B} \supset \bigcup_{i=1}^{\infty}\overline{A_i}$.

Show, by an example, that this inclusion can be proper.
\end{enumerate}
\end{exercise}
\begin{proof}
To prove (a), by mathematical induction, it suffices to prove that $\overline{A \cup B} = \overline{A} \cup \overline{B}$. 
We need to prove that $(A \cup B)'=A' \cup B'$.

Suppose $p \in (A \cup B)'$. 
Assume, for contradiction, that $p$ is neither a limit point of $A$ nor $B$.
Then there are neighborhoods $N_1$ and $N_2$ of $p$ such that $N_1 \cap (A \setminus \{p\})$ and $N_2 \cap (B \setminus \{p\})$ are empty, respectively.

Let $N = N_1 \cap N_2$.
Then we have
\begin{align*}
    [N \cap (A \setminus \{p\})] \cup [N \cap (B \setminus \{p\})] &= N \cap [(A \setminus \{p\}) \cup (B \setminus \{p\})] \\
    &=N \cap [(A \cup B) \setminus \{p\}] = \varnothing.
\end{align*}
This contradicts our assumption. Therefore, $(A \cup B)' \subset A' \cup B'$.

Conversely, suppose $p \in A' \cup B'$.
Then $p \in A'$ or $p \in B'$.
If $p \in A'$, then every neighborhood of $p$ contains a point $q \neq p$ such that $q \in A \subset A \cup B$, hence $p \in (A \cup B)'$.
Similarly, if $p \in B'$ then $p$ is a limit point of $A \cup B$, hence $p \in (A \cup B)'$. 
This proves $A' \cup B' \subset (A \cup B)'$.

We have shown that $(A \cup B)' \subset A' \cup B'$ and $A' \cup B' \subset (A \cup B)'$. 
Hence $(A \cup B)'=A' \cup B'$ and it follows that 
\begin{align*}
    \overline{A \cup B}&=(A \cup B) \cup (A \cup B)'=(A \cup B) \cup (A' \cup B')\\
    &=(A \cup A') \cup (B \cup B') = \overline{A} \cup \overline{B}.
\end{align*}
This is the desired result.

Next, let us prove (b). 
Suppose $p \in \bigcup_{i=1}^{\infty} \overline{A_i}$.
Then $p \in \overline{A_i}$ for some positive integer $i$.
If $p \in A_i$ then $p \in B$.
If $p \in A_i'$ then every neighborhood of $p$ contains a point $q \neq p$ such that $q \in A_i \subset \bigcup_{i=1}^{\infty}A_i=B$.
Therefore, $p$ is a limit point of $B$ and $p \in B'$. 

Hence, if $p \in \bigcup_{i=1}^{\infty} \overline{A_i}$ then $p \in B$ or $p \in B'$. This proves $\bigcup_{i=1}^{\infty}\overline{A_i} \subset \overline{B}$.

Finally, we show that this inclusion can be proper.
For each $n \in \N$, define
\[
A_n=\left\{x \in \R \mid \frac{1}{n}<x \right\}.
\]
Then we see that $B$ is the set of positive real numbers.
Observe that $0$ is not a limit point of $A_n$ and $0 \notin A_n$ for $n=1,2,3, \ldots$.
It follows that $0 \notin \bigcup_{i=1}^{\infty}\overline{A_i}$. But it is clear that $0 \in \overline{B}$.
\end{proof}
\begin{exercise}
Is every point of every open set $E \subset \R^2$ a limit point of $E$?
Answer the same question for closed sets in $\R^2$.
\end{exercise}
\begin{proof}
Suppose $p \in E$ is not a limit point of $E$.
Then there exists a neighborhood $N'$ of $p$ such that $N' \cap E = \{p\}$.
Since $E$ is open, choose a neighborhood $N''$ of $p$ such that $N'' \subset E$.
Then we have 
\[
N= (N'' \cap N') \subset (E \cap N') = \{p\}.
\]
Since $p \in N$, it follows that $N=\{p\}$. Since $N$ is an intersection of two neighborhoods of the same point $p$ in $\R^2$, $N$ is a neighborhood of $p$.
But in $\R^2$, every neighborhood is infinite.
This is a contradiction.

Next, let $F=\{\mathbf{x}\} \subset \R^2$. It is clear that $F$ is a closed subset of $\R^2$ but $\mathbf{x} \in F$ is not a limit point of $F$.
\end{proof}
\begin{exercise}
Let $K \subset \R^1$ consist of $0$ and the numbers $1/n$, for $n=1,2,3,\ldots$. 
Prove that $K$ is compact directly from the definition (without using the Heine-Borel theorem).
\end{exercise}
\begin{proof}
Let $\{G_{\alpha}\}$ be an open cover of $K$.
Since $0 \in K$, there exists an $\alpha_0$ such that $0 \in G_{\alpha_0}$.
Since every $G_\alpha$ is open, choose a neighborhood $N$ of $0$ with radius $r$ such that $N \subset G_{\alpha_0}$.

If $r>1$ then $K \subset G_{\alpha_0}$ hence we are done. 

If $0 < r \leq 1$ then the set 
\begin{equation*}
E=\left\{1/n \mid n \in \N, r \leq 1/n \right\}
\end{equation*}
is finite by the Archimedean Property.
Let $x_1, \ldots, x_k$ be the points of $E$.
Since $x_i \in K \subset \bigcup_{\alpha} G_{\alpha}$ ($1 \leq i \leq k$), there exists $\alpha_i$ such that $x_i \in G_{\alpha_i}$ for $i=1, \ldots, k$.
Thus, we have finitely many indices $\alpha_0, \alpha_1, \ldots, \alpha_k$ such that
\[
K \subset G_{\alpha_0} \cup G_{\alpha_1} \cup \cdots \cup G_{\alpha_k}.
\]

We conclude that $K$ is compact.
\end{proof}

\begin{exercise}
Give an example of an open cover of the segment $(0,1)$ which has no finite subcover.
\end{exercise}
\begin{proof}
For each $n \in \N$, let
\[
G_n=\left\{x \in \R \mid 0<x<1-\frac{1}{n}\right\}.
\]
If $0<x<1$, then there exists a positive integer $n$ such that $1<n(1-x)$, hence $0<x<1-\frac{1}{n}$ so that $x \in G_n$.
It follows that the collection $\{G_n\}_{n \in \N}$ is an open cover of $(0,1)$.
But it is clear that any finite subcollection of $\{G_n\}_{n \in \N}$ cannot cover $(0,1)$.
\end{proof}

\begin{exercise} \,
\begin{enumerate}
\item[(a)] If $A$ and $B$ are disjoint closed sets in some metric space $X$, prove that they are separated.
\item[(b)] Prove the same for disjoint open sets.
\item[(c)] Fix $p \in X$, $\delta > 0$, define $A$ to be the set of all $q \in X$ for which $d(p,q) < \delta$, define $B$ similarly, with $>$ in place of $<$. Prove that $A$ and $B$ are separated.
\item[(d)] Prove that every connected metric space with at least two points is uncountable.

\emph{Hint}: use(c).
\end{enumerate}
\end{exercise}

\begin{proof}\,
\begin{enumerate} 
\item[(a)] Suppose $A$ and $B$ are disjoint closed subsets of a metric space $X$. 
Since $A = \overline{A}$ and $B= \overline{B}$ by the Theorem 2.27, we have $\overline{A} \cap B=A \cap \overline{B}=A \cap B = \varnothing$.
This proves (a).
\item[(b)] Let $A$ and $B$ be disjoint open subsets of $X$.
Since $A \cap B = \varnothing$, we have $A \subset B^c$. 
Since $B^c$ is closed (notice $B$ is open), the Theorem 2.27 implies that $\overline{A} \subset B^c$, hence $\overline{A} \cap B = \varnothing$.
By symmetry, the same argument proves that $A \cap \overline{B} = \varnothing$.
Therefore, $A$ and $B$ are separated.
\item[(c)] Fix $p \in X$ and $\delta>0$.
Then by the definition of $A$, it is a neighborhood in $X$, hence it is open.

We wish to prove that $B$ is open.
Let $q \in B$ be given and put $r=d(p,q)-\delta>0$.
If $x \in X$ and $d(q,x)<r$ then 
\begin{align*}
    d(p,x) & \geq d(p,q) - d(q,x) \\
    &=r+\delta -d(q,x)>\delta,
\end{align*}
hence $x \in B$.
It follows that $N_r(q) \subset B$ and therefore $B$ is open.

Since $A$ and $B$ are disjoint open sets in $X$, (b) shows that they are separated.
\item[(d)] Let $X$ be a connected metric space.
Suppose $p$ and $q$ are two distinct points in $X$ and put $r=d(p,q)>0$.
For every $\delta \in (0,r)$, define 
\begin{align*}
A_{\delta}=\left\{x \in X \mid d(p,x)<\delta\right\}, \qquad B_{\delta}=\left\{x \in X \mid d(p,x)>\delta\right\}.
\end{align*}
It follows from (c) that $A_{\delta}$ and $B_{\delta}$ are separated.
Thus it is impossible that $X=A_{\delta} \cup B_{\delta}$ and hence there exists a point $x_{\delta} \in X$ such that $d(p,x_{\delta})=\delta$ for each $\delta \in (0,r)$.

Let $f(\delta)=x_{\delta}$ for $0<\delta<r$.
If $\delta_1 \neq \delta_2$ then $f(\delta_1) \neq f(\delta_2)$, hence $f$ is a 1-1 function from $(0,r)$ to $X$.
Since $(0,r)$ is an uncountable subset of $\R$, we conclude that $X$ is uncountable.
\end{enumerate}
\end{proof}

